\begin{textsample}{NEG Dim 1 – 2015\_c5 – Score 8.00 – 2015\_c5/enade\_2015\_admin\_p\_0f213d421a.txt}  \label{ex:f1_neg_350}
CAPÍTULO 5 DISTRIBUIÇÃO DOS CONCEITOS 5.1 PANORAMA NACIONAL DA DISTRIBUIÇÃO DOS CONCEITOS O Conceito Enade é um indicador de qualidade que avalia o desempenho dos estudantes a partir dos resultados do Enade.
Os valores possíveis vão de 1 ( pior situação ) a 5 ( melhor situação ).
As Unidades de Observação com apenas um ou sem nenhum concluinte participante não obtêm o Conceito Enade, ficando Sem Conceito ( SC ).
A Tabela 5.1 apresenta a quantidade e distribuição de cursos de Administração Pública participantes do Enade/2015, por faixa de conceito e Grande Região.
Mais uma vez ressaltando que a diferença entre os cursos tabulados neste capítulo e no capítulo 2 corresponde aos cursos sem conceito, em princípio, aqueles sem alunos concluintes que participassem da prova ou que tivessem tão somente um aluno concluinte18.
Observando -se os dados da Tabela 5.1, nota -se que, dos 51 cursos participantes, 20 ( 39,2 \% ) classificaram -se com conceito 3, o valor modal.
Este foi também o conceito modal na região Sul ( 50,0 \% ).
Os conceitos foram distribuídos em duas notas contíguas na região Sudeste ( conceitos 3 e 4, correspondendo a 36,8 \%, \textbf{cada} ) e em duas notas disjuntas na Centro-Oeste ( conceitos 1 e 3, correspondendo a 33,3 \%, \textbf{cada} ).
Na região Nordeste, a moda foi o conceito 2 ( 50,0 \% ).
A região Norte participou com apenas dois cursos, sendo que um curso obteve o conceito 3 e o outro curso ficou sem conceito ( SC ).
Os conceitos 2 e 4 aparecem como os segundos de maior frequência em nível nacional ( 11 cursos, \textbf{cada}, correspondendo a 21,6 \% ) e o conceito 1, o terceiro ( 5,9 \%, correspondendo a três cursos ).
Houve, ainda, dois cursos ( 3,9 \% ) que receberam conceito 5.
Dos 51 cursos de Administração Pública, quatro cursos ( 7,8 \% ) ficaram sem conceito ( SC ). 18 Estes cursos com somente um concluinte não foram considerados no capítulo 2.
A região Norte participou com dois cursos ou 3,9 \% do total nacional.
Nesta região, um curso ( 50,0 \% ) recebeu conceito 3 e um outro ( 50,0 \% ) ficou sem conceito ( SC ).
A região Nordeste participou com 16 cursos ou 31,4 \% do total nacional.
Nesta região, oito cursos, 50,0 \% em termos regionais, obtiveram conceito 2, o conceito modal para a região.
Os conceitos 1 e 4 foram atribuídos a um curso \textbf{cada} ( 6,3 \%, \textbf{cada} ).
Seis cursos ( 37,5 \% ) obtiveram o conceito 3 e nenhum curso obteve o conceito 5 ou ficou sem conceito ( SC ) nessa região.
Dos 19 cursos participantes da região Sudeste ( 37,3 \% do total nacional ), os conceitos 3 e 4 foram distribuídos de forma bimodal, visto que foram obtidos por sete cursos \textbf{cada} ( 36,8 \%, \textbf{cada} ).
O conceito 2 foi atribuído a um curso ( 5,3 \% ) e o conceito 5 a dois cursos ( 10,5 \% ).
Nenhum curso recebeu o conceito 1 e dois cursos ( 10,5 \% ) ficaram sem conceito ( SC ).
Dos oito cursos da região Sul, há predominância do conceito 3 com 50,0 \% ( conceito modal ), ou seja, quatro dos oito cursos participantes na região Sul.
O conceito 2 foi atribuído a um curso ( 12,5 \% ) e o conceito 4, a dois cursos ( 25,0 \% ).
Nenhum curso recebeu conceito 1 ou 5.
Um curso ( 12,5 \% ) na região Sul ficou sem conceito ( SC ).
Já dos seis cursos participantes na região Centro-Oeste, tem -se que os conceitos 1 e 3 foram distribuídos de forma bimodal, visto que foram obtidos por dois cursos \textbf{cada} ( 33,3 \%, \textbf{cada} ).
Nesta região, nenhum curso ficou sem conceito ( SC ) ou com conceito 5.
Dois cursos foram alocados nos conceitos 2 e 4, um curso \textbf{cada} ( 16,7 \%, \textbf{cada} ).
O Gráfico 5.1 apresenta a distribuição do Conceito Enade segundo Grande Região : poligonais mais à esquerda representam Grandes Regiões com pior distribuição, e poligonais mais à direita, Grandes Regiões com melhores conceitos.
As regiões Centro-Oeste ( linha preta ) e Nordeste ( linha verde ) apresentam os cursos em IES com conceitos mais baixos, ao passo que a região Sudeste ( linha vermelha ) apresenta as distribuições com valores maiores.
Cursos em IES das regiões Norte ( linha azul ) e Sul ( linha roxa ) apresentam uma situação intermediária. 5.2 CONCEITOS POR CATEGORIA ADMINISTRATIVA E POR GRANDE REGIÃO A Tabela 5.2 apresenta a distribuição dos cursos participantes do Enade/2015 de Administração Pública, por Categoria Administrativa, de acordo com os conceitos por eles alcançados, segundo as Grandes Regiões brasileiras.
Dos 51 cursos participantes, 47 ( 92,2 \% ) eram ministrados em Instituições Públicas e quatro ( 7,8 \% ), em Instituições Privadas.
De acordo com as informações da Tabela 5.2, em termos nacionais, entre Instituições Públicas e Privadas nota -se que as duas IES que receberam o conceito 5 eram Públicas.
Dos 47 cursos participantes de IES Públicas, o conceito 3 foi o valor modal, atribuído a 20 cursos ( 42,6 \% ).
Entre os demais cursos de Instituições Públicas participantes, três cursos ( 6,4 \% ) foram avaliados com conceito 1, 20 cursos obtiveram conceitos 2 e 4 ( dez cursos \textbf{cada}, 21,3 \% da categoria ) e, como já comentado, dois foram alocados ao conceito 5 ( 4,3 \% ).
Nesta categoria, dois cursos ficaram sem conceito ( SC ).
Na rede privada, dos quatro cursos oferecidos, dois ficaram alocados nos conceitos 2 e 4 ( um curso \textbf{cada} ) e dois ficaram sem conceito ( SC ).
Na análise por região, observa -se que, na região Norte, as Instituições Públicas concentraram todos os dois cursos participantes, sendo que um curso foi avaliado com o conceito 3 e o outro ficou sem conceito ( SC ).
Na região Nordeste, a rede privada contou com apenas um curso participante que recebeu o conceito 2.
As Instituições Públicas dessa região participaram com 15 cursos ( 93,8 \% do total regional ), dos quais sete ( 46,7 \% ) obtiveram conceito 2, o valor modal, seguidos por seis cursos que receberam conceito 3 ( 40,0 \% ).
Dois cursos foram alocados nos conceitos 1 e 4, um curso \textbf{cada}.
Nenhum curso ficou sem conceito ( SC ) ou recebeu conceito 5 nesta combinação de categoria e região.
Na região Sudeste, a proporção de cursos da rede privada, 10,5 \%, corresponde a dois dos 19 cursos participantes.
Nesta categoria, na região Sudeste, um curso recebeu conceito 4, enquanto o outro ficou sem conceito ( SC ).
Entre os 17 cursos oferecidos por Instituições Públicas na região Sudeste, o conceito modal foi 3 ( sete cursos ), e os demais cursos receberam os conceitos 2 ( um curso ), 4 ( seis cursos ) e 5 ( dois cursos ).
Na região Sudeste, nesta categoria, somente um curso ficou sem conceito ( SC ) e nenhum curso recebeu conceito 1.
Na região Sul, as Instituições Privadas foram representadas por apenas um curso que ficou sem conceito ( SC ), 12,5 \% do total regional.
As Instituições Públicas na região Sul participaram com sete cursos ( 87,5 \% ), e o conceito modal foi 3, com quatro cursos.
Um curso foi avaliado com conceito 2 e dois cursos obtiveram conceito 4.
Nenhum curso ficou sem conceito ( SC ) ou recebeu os conceitos 1 ou 5.
Na região Centro-Oeste, as Instituições Públicas concentraram todos os seis cursos participantes, não havendo a participação de cursos de Instituição Privada.
Dos seis cursos, quatro receberam os conceitos 1 e 3, distribuídos de forma bimodal, uma vez que foram obtidos por 2 cursos \textbf{cada} ( 33,3 \%, \textbf{cada} ).
Nesta região, nenhum curso ficou sem conceito ( SC ) ou com conceito 5.
Dois cursos foram alocados nos conceitos 2 e 4, um curso em \textbf{cada} ( 16,7 \%, \textbf{cada} ).
O Gráfico 5.2 apresenta a distribuição do Conceito Enade segundo a Categoria Administrativa da IES.
Os cursos em IES Públicas ( linha azul ) apresentam uma poligonal mais à direita do que os em IES Privadas ( linha verde ) na primeira metade do acumulado e, consequentemente, uma distribuição de Conceitos Enade com valores maiores, enquanto na segunda metade essa situação se inverte. 5.3 CONCEITOS POR ORGANIZAÇÃO ACADÊMICA E POR GRANDE REGIÃO Na Tabela 5.3, encontra -se a distribuição dos conceitos atribuídos aos cursos participantes do Enade/2015 na Área de Administração Pública, por Organização Acadêmica, segundo as Grandes Regiões brasileiras.
Dos 51 cursos de Administração Pública participantes, 46 eram oferecidos em Universidades ; dois cursos em Centro Universitário ; e três cursos em Faculdades.
Esta distribuição corresponde, respectivamente, a 90,2 \%, 3,9 \% e 5,9 \% dos cursos.
De acordo com os dados apresentados, dos dois cursos avaliados com conceito 5, um era vinculado a Universidade e o outro curso a Faculdade.
Dois cursos oferecidos em Universidades ficaram sem conceito ( SC ).
Esse tipo de Organização Acadêmica teve o conceito 3 como modal, com 20 cursos ( 43,5 \% ).
Os demais cursos de Universidades avaliados receberam os conceitos 1 ( três cursos ), 2 e 4 ( dez cursos, \textbf{cada} ), e conceito 5 ( um curso, como já mencionado ).
Os Centros Universitários contaram com apenas dois cursos, que ficaram sem conceito ( SC ).
Nas Faculdades, os três cursos ficaram com conceitos 2, 4 e 5, um curso \textbf{cada}.
Considerando -se separadamente as regiões brasileiras, verifica -se que, na região Norte, as Universidades concentraram a totalidade dos cursos da região.
Dos dois cursos participantes, um curso ficou sem conceito ( SC ) e o outro recebeu conceito 3.
Os Centros Universitários e as Faculdades da região Norte não foram representados por nenhum curso.
Na região Nordeste, as Universidades participaram com 15 dos 16 cursos na Área de Administração Pública.
O conceito modal foi 2, com sete cursos, enquanto o conceito 3 foi atribuído a seis cursos.
Dois cursos obtiveram conceito 1 e 4, um curso \textbf{cada}.
Nenhum dos cursos oferecidos por Universidades no Nordeste ficou sem conceito ( SC ) ou recebeu conceito 5.
Os Centros Universitários não tiverem nenhum curso participante na região Nordeste.
As Faculdades foram representadas por apenas um curso, que recebeu o conceito 2.
Na região Sudeste, as Universidades concentraram 16 dos 19 cursos da região.
Entre os cursos em Universidades, nessa região, o conceito modal foi 3 com sete cursos, enquanto um curso ficou sem conceito ( SC ) e nenhum recebeu o conceito 1.
Os demais cursos receberam os conceitos 4 ( seis cursos ), e dois cursos obtiveram os conceitos 2 e 5 ( um curso \textbf{cada} ).
Os Centros Universitários tiveram apenas um curso participante na região Sudeste, que ficou sem conceito ( SC ).
As Faculdades foram representadas por dois cursos na região Sudeste que receberam os conceitos 4 e 5 ( um curso \textbf{cada} ).
Nenhum curso ficou sem conceito ( SC ) nessa categoria e região.
Dos oito cursos da região Sul, sete eram de Universidades, para os quais o conceito modal foi 3, com quatro cursos.
Nesse tipo de organização, nenhum dos cursos ficou sem conceito ( SC ) ou recebeu os conceitos 1 ou 5.
Os demais receberam os conceitos 2 ( um curso ) e 4 ( dois cursos ).
O único curso participante da categoria Centro Universitário da região Sul ficou sem conceito ( SC ).
As Faculdades na região Sul não contaram com nenhum curso participante na região.
Na região Centro-Oeste, verifica -se que as Universidades concentraram a totalidade dos seis cursos da região.
Entre eles, quatro cursos obtiveram conceitos 1 e 3 ( dois cursos \textbf{cada} ) e dois cursos foram avaliados com os conceitos 2 e 4 ( um curso \textbf{cada} ).
Nenhum curso ficou sem conceito ( SC ) ou recebeu conceito 5.
Os Centros Universitários e as Faculdades da região Centro-Oeste não foram representados por nenhum curso.
O Gráfico 5.3 apresenta a distribuição do Conceito Enade segundo a Organização Acadêmica da IES.
Os cursos em Faculdades ( linha vermelha ) apresentam uma poligonal mais à direita do que a dos cursos em Universidade, denotando uma melhor distribuição de Conceitos Enade.
No outro extremo, tem -se os cursos oferecidos em Universidades ( linha azul ) com a poligonal mais à esquerda e, consequentemente, uma distribuição de Conceitos Enade com menores valores.
No entanto, há que se considerar que as Faculdades se fizeram representar por, apenas, três cursos e as Universidades por 46 cursos.
IES de Centros Universitários ( linha verde ) não foram representados no Gráfico, visto que os dois únicos cursos desta categoria ficaram sem conceito ( SC ).

% matched lemmas: cada
\end{textsample}
